% titel: Blatt 8, A4: Phasendiagramm eines reinen Stoffs lesen
% skript: Mortimer Kap. 11 (Phasendiagramme)

% Teil I: Werkzeuge, kein Lösungsweg.

Ein Phasendiagramm zeigt für einen reinen Stoff, welche Phase bei gegebenem
Druck $p$ und gegebener Temperatur $T$ stabil ist. Es ist die zentrale
Abbildung, um Übergänge fest, flüssig, gasförmig zu beschreiben.

\begin{werkzeugkasten}[Phasendiagramm lesen]
\begin{itemize}
  \item \textbf{Regionen:} fest (s), flüssig (l), gasförmig (g) --- jeweils in
        einem Bereich der $p$-$T$-Fläche.
  \item \textbf{Koexistenzlinien:} auf der Linie stehen zwei Phasen im
        Gleichgewicht; beim Überqueren wechselt der Aggregatzustand.
  \item \textbf{Tripelpunkt (`Tripelpunkt'):} Punkt, an dem alle drei Linien
        zusammentreffen --- f, l, g sind gleichzeitig im Gleichgewicht.
  \item \textbf{kritischer Punkt (`kritischer Punkt'):} Ende der
        Dampfdruckkurve; oberhalb davon gibt es keine $p,T$-unterscheidbare
        Gas/Flüssig-Grenze mehr.
  \item Typisches Wasser markiert: die Erhöhung des Schmelzpunkts hängt von
        der Steigung der fest/flüssig-Linie ab.
\end{itemize}
\textit{Skript:} Mortimer Kap. 11 (Phasendiagramme, kritischer Punkt).
\end{werkzeugkasten}

\begin{achtung}[Darauf achten]
\begin{itemize}
  \item Ein Punkt \emph{in} einer Region steht für genau eine Phase.
  \item Die Punkte B und D sind die beiden ausgezeichneten Punkte des
        Diagramms (Tripelpunkt, kritischer Punkt).
  \item Auf einer Koexistenzlinie laufen die Umwandlungen in beide
        Richtungen: schmelzen/gefrieren bzw. sieden/kondensieren.
  \item Die Steigung der fest/flüssig-Linie bestimmt, ob der Schmelzpunkt
        mit dem Druck steigt oder fällt (beim Wasser steigt die Linie nach
        links, d.~h. der Schmelzpunkt sinkt bei Druckerhöhung).
\end{itemize}
\end{achtung}
